\documentclass[12pt]{article}

% ============================================================
% PACKAGES
% ============================================================

\usepackage{amsmath}      % Mathematical environments: align, equation, etc.
\usepackage{amsthm}       % Theorem, definition, proof environments
\usepackage{amsfonts}     % Mathematical fonts such as \mathbb
\usepackage{amssymb}      % Additional mathematical symbols
\usepackage{mathtools}    % Extensions to amsmath
\usepackage{geometry}     % Page margins
\usepackage{enumitem}     % Customize enumerate/itemize
\usepackage{mdframed}     % Framed boxes
\usepackage{tikz}         % Draw diagrams
\usepackage{hyperref}     % Clickable references and links

\geometry{
    letterpaper,
    margin=0.8in
}

\linespread{1.2}


% ============================================================
% NUMBER SYSTEMS
% ============================================================
% Instead of writing \mathbb{R} every time, use \RR.

\newcommand{\RR}{\mathbb{R}}   % Real numbers
\newcommand{\CC}{\mathbb{C}}   % Complex numbers
\newcommand{\NN}{\mathbb{N}}   % Natural numbers
\newcommand{\ZZ}{\mathbb{Z}}   % Integers
\newcommand{\QQ}{\mathbb{Q}}   % Rational numbers


% ============================================================
% COMMON MATHEMATICAL NOTATION
% ============================================================
% These commands are useful for frequently used notation.

\newcommand{\set}[1]{\left\{#1\right\}}
% Example:
%   \set{x \in \RR : x > 0}
% produces:
%   { x in R : x > 0 }

\newcommand{\abs}[1]{\left|#1\right|}
% Example:
%   \abs{x-y}
% produces:
%   |x-y|

\newcommand{\norm}[1]{\left\|#1\right\|}
% Example:
%   \norm{x}
% produces:
%   ||x||

\newcommand{\inner}[2]{\left\langle #1,#2\right\rangle}
% Example:
%   \inner{x}{y}
% produces:
%   <x,y>

\newcommand{\parens}[1]{\left(#1\right)}
% Example:
%   \parens{x+y}

\newcommand{\brac}[1]{\left[#1\right]}
% Example:
%   \brac{a,b}

\newcommand{\angles}[1]{\left\langle #1\right\rangle}
% Example:
%   \angles{x,y}


% ============================================================
% THEOREM / THEORY COMPONENTS
% ============================================================

% "definition" style is upright text.
% Use this for definitions, examples, and remarks.

\theoremstyle{definition}

% ------------------------------------------------------------
% DEFINITION
% ------------------------------------------------------------
% Use when introducing a formal mathematical definition.
%
% Usage:
%
% \begin{definition}[Metric]
% A metric on a set X is a function ...
% \end{definition}
%
% [Metric] is optional and gives the definition a name.

\newtheorem{definition}{Definition}[section]


% ------------------------------------------------------------
% EXAMPLE
% ------------------------------------------------------------
% Use to give a concrete example of a definition/theorem.
%
% Usage:
%
% \begin{example}
% Consider X = \RR with the usual distance.
% \end{example}

\newtheorem{example}[definition]{Example}


% ------------------------------------------------------------
% REMARK
% ------------------------------------------------------------
% Use for observations, comments, or small additional facts.
% It is NOT intended for major mathematical results.
%
% Usage:
%
% \begin{remark}
% This definition does not require X to be a subset of \RR^n.
% \end{remark}

\newtheorem{remark}[definition]{Remark}


% ============================================================
% THEOREM-LIKE COMPONENTS
% ============================================================

% "plain" style makes the mathematical statements italic.

\theoremstyle{plain}


% ------------------------------------------------------------
% THEOREM
% ------------------------------------------------------------
% Use for an important/general mathematical result.
%
% Usage:
%
% \begin{theorem}[Cauchy--Schwarz Inequality]
% For all x,y \in \RR^n,
% \[
% |\inner{x}{y}| \leq \norm{x}\norm{y}.
% \]
% \end{theorem}

\newtheorem{theorem}[definition]{Theorem}


% ------------------------------------------------------------
% PROPOSITION
% ------------------------------------------------------------
% Use for a useful mathematical result that is usually less
% central than a theorem.
%
% Usage:
%
% \begin{proposition}
% Every finite-dimensional subspace of \RR^n is closed.
% \end{proposition}

\newtheorem{proposition}[definition]{Proposition}


% ------------------------------------------------------------
% LEMMA
% ------------------------------------------------------------
% Use for an auxiliary result whose main purpose is to help
% prove another theorem/proposition.
%
% Usage:
%
% \begin{lemma}
% ...
% \end{lemma}

\newtheorem{lemma}[definition]{Lemma}




% ------------------------------------------------------------
% COROLLARY
% ------------------------------------------------------------
% Use for a result that follows relatively directly from a
% theorem/proposition that was just established.
%
% Usage:
%
% \begin{corollary}
% ...
% \end{corollary}

\newtheorem{corollary}[definition]{Corollary}


% ============================================================
% PROOF
% ============================================================
% The proof environment comes from amsthm.
%
% Usage:
%
% \begin{proof}
% Suppose that ...
%
% Therefore, ...
% \end{proof}
%
% LaTeX automatically adds the QED symbol at the end.


% ============================================================
% CUSTOM BOX: IDEA
% ============================================================
% Use this when you want to record intuition or the main idea
% behind a mathematical concept.
%
% Usage:
%
% \begin{idea}
% \textbf{Key idea.}
% A metric gives us a notion of distance...
% \end{idea}

\newmdenv[
    linewidth=1pt,
    skipabove=10pt,
    skipbelow=10pt
]{idea}


% ============================================================
% CUSTOM BOX: WARNING
% ============================================================
% Use this for common mistakes, distinctions, or things that
% are easy to confuse.
%
% Usage:
%
% \begin{warning}
% Do not confuse continuity with uniform continuity.
% \end{warning}

\newmdenv[
    linewidth=1pt,
    skipabove=10pt,
    skipbelow=10pt
]{warning}

% ============================================================
% CUSTOM BOX: NOTE
% ============================================================
% Use this for a short personal note, useful observation,
% connection between concepts, or explanation.
%
% Unlike "Remark", this is intended for informal notes.
%
% Usage:
%
% \begin{note}
% A sequence can be viewed as a function from $\NN$ to $\RR$.
% \end{note}

\newmdenv[
    linewidth=1pt,
    skipabove=10pt,
    skipbelow=10pt
]{note}


% ============================================================
% DOCUMENT
% ============================================================

\begin{document}


% ============================================================
% TITLE
% ============================================================

\begin{center}
    {\Large \textbf{Least Upper Bound Property}}\\[0.5em]
    {\large Lecture 4}
\end{center}

% Automatically generate the table of contents.
\tableofcontents

\newpage
\section{L.U.B. property in the Reals}
\begin{note}
In the last note, we know that, $\RR$ has the Least Upper Bound property (which mean any non-empty subset of $\RR$ if it is bounded, then it has least upper bound, which still belongs to $\RR$). So, the natural question come up with you right now (maybe): What is that least upper bound?
\end{note}

\begin{proposition}
Given a subset $A$ of $\RR$, which is a collection of cuts in $\RR$, which is all upper bounded by $\beta$ (not necessarily the least upper bound). Then, $\gamma = \bigcup\{\alpha \in A\}$ is a cut, and $\sup A = \gamma$. 
\end{proposition}


\begin{note}
At this point I can just say that, a cut $\alpha$ represents the $\sup$ of its set of rationals in $\RR$ (This is enough to imagine adding 2 real numbers or something else,... the result of adding reals is just a cut, but in practice, we just write all of them in the form of the $\sup$ of cuts. Additionally, you can imagine adding 2 reals is just adding 2 cuts, then take the $\sup$-representation as the result).
\end{note}

\section{Embed the Rationals into the Reals}
\begin{note}
As you know, a cut represent the "$\sup$", so if the $\sup$ is a rational number, then we embedded $\QQ$ into $\RR$.
\end{note}

\begin{definition}[Rational cuts]
Given $q \in \QQ$, then $q^* = \{r\in\QQ:r<q\}$ is a rational cut in $\RR$, which represents a rational number: $q$. 
\end{definition}

\begin{note}
But the problem of irrational numbers is still there. How exactly a bunch of rationals can form a irrational number. Of course you may already recognized it, the key object is the $\sup$ of that set of rationals. In fact, the $\sup$ can be irrational, that's how we obtained the irrationals.
\end{note}

\begin{example}
In $\RR$, there exists a cut to represent $\sqrt{2}$, which is $\{r \in \QQ: r^2<2 \text{ or } r<0\}$. You can see that, the $\sup$ of the given set, is indeed, $\sqrt{2}$.
\end{example}

\begin{definition}[General roots]
We can have 2 ways to define general roots:
\begin{itemize}
    \item \textbf{Definition with $\RR$:} $a^{\frac{1}{n}} = \sup\{r\in \RR: r^n<a\}$.

    \item \textbf{Definition with $\QQ$:} $a^{\frac{1}{n}}=\{r \in \QQ: r^{n}<a \text{ or } r<0\}$
\end{itemize}

\end{definition}

\section{Greatest Lower Bound}
\begin{definition}[Infimum]
The Infimum of a given lower bounded set $E$, denoted as $\inf E$, is the Greatest Lower Bound of $E$.
\end{definition}

\begin{proposition}
$\inf E = -\sup(-E)$
\end{proposition}

\section{More on the Lowest Upper Bound property}

\begin{note}
Now you can recover the imagination for $\RR$ as normal (numbers, not cuts). But the proofs for those propositions still needed cuts. But in the imagination of real numbers, those stuffs are quite easy to understand, or even, trivial. I won't mention the proof in this note.
\end{note}

\begin{proposition}[Archimedean]
$x,y \in \RR, x>0 \Rightarrow \exists n>0, n\in\ZZ: nx>y$.
\end{proposition}

\begin{note}
One of useful applications of Archimedean is that the sequence $nx$ is unbounded.
\end{note}

\begin{proposition}[$\QQ$ is dense in $\RR$]
For any $x,y\in\RR, x<y \text{ then }\exists q\in\QQ: x<q<y$.
\end{proposition}

\begin{proposition}[Some properties of $\sup$]
Here are some properties for the supremum:
\begin{itemize}
    \item $\gamma$ is an upper bound (u.b.) for $A$ if and only if $\sup A \le \gamma$.
    \item If $\forall a \in A$ we have $a\le\gamma$, then $\sup A \le \gamma$.
    \item If $\forall a \in A$ we have $a<\gamma$, then $\sup A \le \gamma$.
    \item To show $\gamma$ is the supremum of a given set, show that $\gamma$ is an u.b., and any $x < \gamma$ is not an u.b. or you can show that $\gamma$ is an u.b. and every u.b. is greater or equal to $\gamma$.
    \item If $A\subset B$, then $\sup A \le \sup B$.
    \item $A+B=\{a+b: a\in A, b\in B\}$ then $\sup(A+B) = \sup A + \sup B$.
    \item $A \cdot B = \{ab: a \in A, b\in B\}$ then if $A, B \in \RR_+$ then $\sup AB = \sup A \cdot \sup B$.
\end{itemize}
\end{proposition}

\end{document}